In order to further research on LM-based systems that can assist scientific progress, we introduce \model and \data, which can help navigate the complex, ever-growing task of scientific literature review. 
\model, a retrieval-augmented system, leverages open-checkpoint LLMs and trained retrieval models to iteratively refine scientific output, addressing challenges such as hallucinations and citation accuracy. 
\data, a novel large-scale benchmark, provides a standardized way to evaluate literature review automation across multiple scientific domains. 
In evaluations using \data, \model demonstrates substantial improvements, outperforming existing systems, including GPT-4o and the concurrent proprietary system PaperQA2. Our expert evaluation across three scientific disciplines reveals that \data, when paired with fully open-checkpoint models and open-access data stores, generates answers that are more helpful than those produced by expert annotators, who required an hour per annotation. This approach also significantly increases coverage. 
\model using our trained 8B and GPT4o  achieves a 51\% and 70\% win rate against human-generated answers. 
We open-source the \model code, data, model checkpoints, datastores, and \data, along with a public demo, to support and accelerate future research efforts. 

\section*{Limitations}
We highlight several limitations of our work in this section. It is important to note that we do not claim that LM-based systems can fully automate scientific literature synthesis. 
To further advance research in this area, we are releasing both \data and \model to the community. 

\paragraph{Limitations of \data.}
There are several limitations to \data. First, due to the cost and time required to engage expert annotators—individuals with either a Ph.D. or are currently pursuing one in relevant fields—the evaluation dataset with human-written answers is relatively small (e.g., 110 for CS-LFQA and 108 for expert-written answers). This limited dataset may introduce statistical variance and potential biases stemming from the specific expertise of the annotators. To support future research in expanding the size and scope of \data, we open-source our data and annotation pipelines.

Second, our automatic evaluation pipelines may not always perfectly capture the quality of generated content. For example, in \textsc{Scholar-CS}, we combine various components (e.g., length, excerpts, rubric items) using heuristically determined weight terms. {Further, we discovered that often annotators asked for specific kinds of ancillary information in their rubrics---background, elaborations and challenges---even though these aspects might not be strictly required to answer the question. In our experiments, we found that LLMs are proficient at generating the background aspects, which can give them an advantage over systems that directly answer a query but do not satisfy all the constraints of the rubrics. Moreover, future systems could potentially take advantage of the stylistic biases in the rubrics and be prompted to address more rubric elements in a way that does not improve answer quality.} Although we carefully analyzed the correlation between final scores and human expert evaluations, there is still room for improvement in refining which aspects should be emphasized and how these scores should be aggregated. Additionally, our evaluations of citation precision and recall are conducted at the sentence level, but we found that some sentences without direct citations are often supported by citations in adjacent sentences. As a result, our precision and recall metrics might be overly strict, potentially underestimating the true citation accuracy.  We also note that our annotations were captured at particular times (July 2024 for \textsc{Scholar-CS} and September 2024 for \textsc{Scholar-Multi}), and may not reflect subsequent scientific developments. 
Researchers who use our evaluation benchmark should ignore papers published after these dates for a fair comparison.

Lastly, \data primarily focuses on computer science, biomedicine, and physics, with no instances from social sciences or other engineering and scientific disciplines. We recognize that our findings may not fully generalize to other domains, particularly those with more restricted access to paper data.

\paragraph{Limitations of \model.}
While \model demonstrates strong performance on \data and in human evaluations, as discussed in the relevant sections, our expert annotators identified several limitations. 
Despite these issues, we believe \model remains a valuable tool for supporting human experts. 

First, as highlighted by our expert annotators, \model does not consistently retrieve the most representative or relevant papers for certain queries. Enhancing retrieval methodologies by incorporating additional information, such as citation networks or metadata like publication recency, could significantly improve its performance. 
\model outputs may contain factual inaccuracies or unsupported information, particularly in versions based on our 8B model, which has limited capacity for instruction-following and scientific knowledge.  
Future work can explore training that further improve \model-8B. 
In parallel, although it is competitive, \model-GPT4o relies on the proprietary GPT4o API, which may evolve over time, making exact reproduction of results challenging. 
\model does not use license-protected papers during inference time. 
There are ongoing discussions on how to ensure fair data use in retrieval-augmented LMs, and we leave the exploration of properly incorporating copyright-protected content to future work. 

We encourage future research to address these limitations and continue improving LM-based systems for scientific literature review. 



\paragraph{Limitations of our human evaluation process.}
In our human evaluations, annotators performed fine-grained assessments on aspects such as coverage, relevance, organization, and usefulness, while other factors such as citation precision and recall were separately evaluated. As a result, when assessing usefulness or pairwise preferences, annotators may have focused more on the overall quality of writing instead of carefully evaluating factual correctness or citation accuracy. We leave more detailed human analysis on citation accuracy, validity, and factuality for future work.

Our evaluations were conducted by 16 Ph.D. students and postdoctoral professionals, and we made an effort to align their expertise with the topics being evaluated. However, since research often demands deep domain knowledge, the annotators may not have captured more nuanced differences for questions outside their immediate areas of expertise. 
Additionally, these evaluations were based on 108 questions spanning three scientific disciplines, meaning that findings may not fully generalize to other fields or domains. 